# Chapter 4 main body

\section{Introduction}\label{ch4-introduction}

\subsection{Motivation}\label{ch4-motivation}
-Motivate the need for fast thermodynamic forward models in geodynamics. Composition space is highly unexplored
-Review current tools (Perple\_X, THERMOCALC, MAGEMin) and their computational bottlenecks. (research and citations)

\subsection{Previous Work}\label{ch4-prev-work}
-Review the state of surrogate modeling in geodynamics (research and citations) (qian and ETH group)

\subsection{Properties and Innovations of nGibbs}\label{ch4-innovations}
-High level introduction to nGibbs and what its innovation is (full thermodynamic emulator, obeys mass balance, differentiable, flexible, central tool for many geology-relevant tools, which are all written in their own languages and have independent learning curves)

\section{Methods}\label{ch4-methods}

\subsection{Network architecture}
\label{sec:continuous}

\subsubsection{Motivation: the discontinuity problem}
\label{sec:continuous-motivation}

Before the introduction of the continuous-saturation framework described below, the early \emph{nGibbs} architecture predicted phase presence with a binary classifier
and phase abundance with a distinct regressor, multiplying the two so that the predicted
abundance is the regressed magnitude gated by a thresholded presence probability. 
This gated model was simple and easy to train, but it has a structural defect.
Because the gate and the magnitude were distinct functions,
nothing tied the magnitude to zero as the gate closes. Therefore, when the presence probability of a phase
crosses its threshold, the predicted abundance jumps discontinuously from a finite
value to exactly zero, which is (for most phases) an unnexpected behavior for a mixed system.

Natural systems of interest to geologists and planetary scientists are typically complex, chemically-mixed systems. These systems, and the thermodyanamic models that describe them, 
behave continuously. In addition to the continuous quality of phase abundance, thermodynamic models like \textsc{MELTS} and \textsc{HeFESTo} also track each phases' chemical affinity, 
which is a continuous quality that measures the free energy cost of a metastable phase's formation introduction to the system. 
Even when the phase is absent, the chemical affinity is a meaningful, continuous quantity that smoothly goes to zero as the phase approaches saturation. 
Likewise, at a phase-out boundary, the abundance of the departing phase goes to zero smoothly, linearly in reaction progress. A gated model is
therefore structurally wrong precisely at the phase boundaries that constitute the most
interesting content of a phase diagram.

The architecture used for all released models removes this defect by making the gate
and the magnitude the same variable, analogously to a prediction of chemical affinity. Some emulators predict a scaled chemical affinity, while 
others simply activate their continuous molar output with a ReLU function. The models trained with a scaled affinity perform slightly better, 
but require a more complex data processing pipeline that cannot be applied to datasets where the data are not structured in paths that hold 
a state variable constant. Because the HeFESTo data were generated along mock-isentropic pressure-temperature paths (and not true isentropes), 
the HeFESTo emulators were trained with the simpler, ReLU activated architecture, hereafter referred to as the affinity-naive architecture. 
The MELTS emulators, which were trained on data generated by isobaric cooling, were trained with scaled affinities with an architecture 
hereafter referred to as affinity-aware.

%=====================================================================================
\subsubsection{Variables}
%=====================================================================================

Here a table is provided to name the key variables used in the equations below that describe the training and function of nGibbs. 
These equations rely additionally on general letter codes to describe the dimensionality of the datathat are used internally within 
the nGibbs codebase and are borrowed here for clairity: P phases, VP compositionally variable phases, C components (endmembers of all phases),
VC components of the compositionally variable phases, E elements and O oxides. Two further sizes are
used here: F, the number of physical non-compositional input features (e.g.\ pressure and temperature, with or without \emphfO\textsubscript{2}),
and B, the number of rows in a mini-batch (the amount of assemblages processed in a single pass during training). 
Every quantity in Table~\ref{tab:key} is per row; a batch stacks B rows. For emulator outputs, the hatted symbol is the prediction and
the plain symbol is the ground truth taken directly from the thermodynamic model (e.g. MELTS or HeFESTo).

\begin{table}[!tbp]
\centering
\small
\renewcommand{\arraystretch}{1.25}
\begin{tabular}{@{}c c c >{\raggedright\arraybackslash}p{10.2cm}@{}}
\toprule
Prediction & Truth & Dim. & Meaning \\
\midrule
\multicolumn{4}{@{}l}{\emph{Inputs}}\\
 & $u$ & F & Physical state (pressure, temperature, \dots), min--max normalised to $[0,1]$ \\
$\hat b$ & $b$ & E & Bulk composition in element moles, normalised to sum to 1. $\hat b$ is the bulk reconstruction $\hat b = \hat m A$ \\
\midrule
\multicolumn{4}{@{}l}{\emph{Emulator outputs}}\\
$\hat n$ & $n$ & P & Phase abundance (moles); 0 for an absent phase \\
$\hat x$ & $x$ & VC & Endmember mole fractions of each variable phase (sum to 1 within a phase: $\sum x_{i} = s_{i} |1-b|$) \\
$\hat m$ & $m$ & C & Extensive endmember components/species in moles;  $\hat m = (\hat x V + b) \circ (\hat n M)$ \\
$\hat s$ & $s$ & P & Phase presence, 1 if present and 0 if absent; $ s = [ n > 0]$ \\
\midrule
\multicolumn{4}{@{}l}{\emph{Quantities used only in training}}\\
$g$ & -- & P & Phase affinity/abundance: the raw output of each phase's mole head; $\hat n = \max(g,0)$ \\
$\tilde n$ & -- & P & Smoothed, positive, non-zero abundance used to evaluate mole loss; $\tilde n(g,T)$: Eq.~\eqref{eq:smooth} \\
-- & $\bar s$ & VC & Phase Presence $s$ projected to phases' endmembers (chemistry mask) $\bar s = s\,M\,V^{\top} \;\in\; \{0,1\}$ \\
\midrule
\multicolumn{4}{@{}l}{\emph{Training constants}}\\
\multicolumn{2}{c}{$T_0$} & P & Reference abundance scale per phase (5th percentile of present abundances) \\
\multicolumn{2}{c}{$T$} & P & Annealed smoothing temperature; $T = a\,T_0$ \\
\multicolumn{2}{c}{$a$} & 1 & Annealing factor, decreased linearly over an episode \\
\multicolumn{2}{c}{$w$} & P & Phase weights \\
\multicolumn{2}{c}{$\omega$} & VC & Endmember weights; $\omega = w\,M\,V^{\top}$ \\
\multicolumn{2}{c}{$\lambda_s,\lambda_x,\lambda_n,\lambda_b$} & 1 & Weights for the binary phase presence, chemistry, molar, and bulk loss terms \\
\midrule
\multicolumn{4}{@{}l}{\emph{Fixed stoichiometric/projection matrices}}\\
\multicolumn{2}{c}{$A$} & C$\times$E & Element moles per component mole \\
\multicolumn{2}{c}{$M$} & P$\times$C & Which components belong to which phase; $M \;\in\; \{0,1\}$ \\
\multicolumn{2}{c}{$V$} & VC$\times$C & Indexing matrix that places the VC components into the full component list; $M \;\in\; \{0,1\}$ \\
\multicolumn{2}{c}{$b$} & C & Which components belong to fixed-composition (pure) phases; $b \;\in\; \{0,1\}$ \\
\bottomrule
\end{tabular}
\caption{Key to the symbols. ``--'' marks a quantity with no prediction or no ground-truth
counterpart.}
\label{tab:key}
\end{table}


\subsubsection{Inputs}

Each emulator takes pressure and one further intensive variable (temperature
for the isothermal models, specific entropy for the isentropic models) and a bulk
composition as inputs. The emulators are trained on and expect the same units for pressure, temperature, and entropy that
the real thermodynamic MELTS and HeFESTo use: Kelvin and GPa for HeFESTo, and Celcius and bars for MELTS. Both expect specific 
entropy in units of J/g/K. nGibbs recieves as input normalized cation atomic abundances, following  $\mathbf{b}\in\mathbb{R}^{E}$ with $\sum_e b_e = 1$. 
Compositions are supplied by the user either as oxide weight percents or as element
molar abundances for any model; oxide inputs are automatically converted internally to the normalized cation abundances expected by nGibbs.

Entropy is offered as a useful alternative to temperature because isentropic paths are the
natural coordinate for the convecting mantle \citep{Davies1977,Brown1981,Asimow1997,Asimow2001,Ganguly2005,Tirone2016}. Supplying $S$ directly removes an inner iteration for $T$ at every pressure step, which is the dominant cost of constructing
an adiabat by conventional means.

The full non-compositional feature vector (i.e. pressure, temperature, entropy, \emph{f}O\textsubscript{2}) is min--max normalized to $[0,1]$. The minima and maxima used for these normalizations are fit on the
training data, and are carried by the emulator and automatically applied during inference.
This avoids presenting the network with inputs on a
different scale than the one it was trained on, leading to errors that would be difficult to diagnose.

The neural networks driving \emph{nGibbs} expect that the input composition vector sum to a total of one. 
During inference, this is explicitly enforced by a renormalization so that, similarly to \textsc{MELTS} and \textsc{HeFESTo}, 
the user can supply a composition that sums to values other than one for convenience. 
The outputs of \emph{nGIbbs}, however, are not rescaled and will always map to a cation sum of one regardless of user input.


\subsubsection{Compositional basis}

Three compositional bases are used. \emph{Components} ($m$) are the endmembers of the reference code's
solution models ($C$ in total across all phases); \emph{oxides} ($O$) are accepted as an input
because it is conventional to measure and report rock compositions as oxides; and \emph{elements} ($E$) are the conserved
quantities that \emph{nGibbs} conserves. Ferric and ferrous iron are distinguished as oxides but not as
elements for models with oxygen buffering, for which $O = E+1$, but closed system models distinguish ferric and ferrous 
iron in the element basis as well so that they are conserved: $O = E$.  Conversion between bases is exact stoichiometry, expressed as
fixed matrices that relate components to oxides and oxides to elements. The product of the two, $A$, is the most commonly
used stoichiometry matrix because it relates model outputs to the bulk elemental composition it conserves. These matrices are model constants,
not fitted parameters.

\emph{nGibbs} models closed systems that conserve oxygen content, as well as open systems where oxygen fugacity is imposed. 
Iron is the major redox-active element in \emph{nGibbs}. Every element besides iron is tracked at a single, 
fixed oxidation state (with a limited exception for carbon) and therefore the oxygen content of the system is fully specified by the ferric-to-total iron ratio. 
The only other redox-active element is carbon in \textsc{MELTS} 1.2 emulators, wherein carbon exists as an oxide in melt, 
fluid, and calcite, and as reduced carbon in graphite. Graphite is represented as C(FeO)\textsubscript{4}(Fe\textsubscript{2}O\textsubscript{3})\textsubscript{-2} in oxide space 
in order to conserve oxygen, which introduces strong coupling between Fe$^{2+}$/Fe$^{3+}$ and graphite for closed system MELTS 120 models. This coupling vanishes for 
open system models where oxygen is not conserved. 
Similarly, native iron phases are represented as (FeO)\textsubscript{3}(Fe\textsubscript{2}O\textsubscript{3})\textsubscript{-1} in oxide space in \emph{nGibbs}.

\emph{nGibbs} accepts one of two inputs that define the redox state: an explicit, distinct definition of both 
FeO and Fe$_2$O$_3$, which the closed-system models conserve; or an oxygen fugacity (\emph{f}O\textsubscript{2}) 
in log unit offset from the fayalite-magnetite-quartz buffer ($\Delta$FMQ), which the open-system models impose. 
Oxygen buffering is neither available in \textsc{HeFESTo} or its \emph{nGibbs} emulator. 
In these open models, only cation mass is conserved and the oxygen contents of the system is a consequence 
of the predicted phase assemblage, except for the liquid, for which total iron is speciated following the procedure 
of \citet{kress1991}.

\subsubsection{Forward pass}

The network sees the input vector state variables $u$ and bulk composition $b$, represented as the vector $[u, b]$ and:
\begin{enumerate}[nosep]
  \item a shared encoder produces a central working representation of the system's state;
  \item one small network per phase gives the signed abundance/affinity $g$;
  \item one softmax-activated network per chemically-variable phase gives that phase's endmember fractions $\hat x$.
\end{enumerate}

During inference, the predicted abundance is the rectified affinity,
  \begin{equation}
    \hat n = \max(g, 0).
  \end{equation}
\end{enumerate}
Because $\max(g,0)$ is continuous at $g=0$, a phase leaving the assemblage is a kink in
$\hat n$ rather than a jump. Additionally, phases and chemical endmembers that contain an element that is absent from $b$ are forced to be zero. 

Component moles and the reconstructed bulk composition follow from stoichiometry:
\begin{align}
  \hat m &= (\hat x V + b) \circ (\hat n M), \label{eq:m}\\
  \hat b &= \hat m A. \label{eq:bhat}
\end{align}
Here $\circ$ is the element-wise product, so $\hat m$ is each component's phase abundance
multiplied by its fraction in that phase. Mass balance (adjusting $\hat m$ until $\hat b = b$) is enforced only at inference (\S\ref{sec:massbalance}).  

\subsubsection{Thermodynamic justification for softmax activated phase composition prediction}\label{ch4-softmax}

The softmax activation function produces a positive, nonzero output that sums to 1, and is most often applied in machine learning for classification problems in tandem with 
cross-entropy loss, as its outputs are often interpreted as a probability or liklihood distribution over categorical classes \citep{holm2023,pearce2021,li2019,yang2018}.  
For class $i$ of $\phi$ categorical classes, the softmax function is $\frac{e^{z_i}}{\sum_{j\in\phi} e^{z_j}}$ where $z_i$ is the logit (linear dense layer output) for class $i$. 

The softmax function is used here for a bounded regression problem, to predict the mole fractions of chemical endmembers in a phase, which are also positive, nonzero, and sum to 1. 
Originally, we chose the softmax function because of this quality, and early tests showed it outperformed alternative activations (such as a renormalization of ReLU outputs) 
for predicting chemical endmember fraction. Although this application of the softmax is unusual, it is justified in that its functional form mirrors that of the thermodynamic 
problem of predicting chemical endmember fractions in a phase at equilibrium.

For a general solution, the chemical potential of endmember $i$ in phase $\phi$ is
$\mu_i = \mu_i^{\circ}(P,T) + RT\ln a_i$, with activity $a_i = \gamma_i x_i$ (e.g.\
Anderson, 2005; Ganguly, 2008). The activity coefficient $\gamma_i \equiv a_i/x_i$
collects the departure from ideal molecular mixing, including excess (interaction)
terms and multi-site configurational effects. At equilibrium, $\mu_i$ equals the apparent chemical potential
$\mu_i^{*}$ that the coexisting assemblage imposes on that endmember, so
\begin{equation}
  \ln x_i \;=\; z_i \;\equiv\; \frac{\mu_i^{*} - \mu_i^{\circ}(P,T) - RT\ln\gamma_i}{RT}.
  \label{eq:nonideal_z}
\end{equation}

The standard state chemical potential $\mu_i^{\circ}(P,T)$ and the unideal contribution ($RT\ln\gamma_i$) may be gathered together in 
an effective reference chemical potential $\tilde{\mu}_i^{\circ}(P,T) \equiv \mu_i^{\circ}(P,T) + RT\ln\gamma_i$.

\begin{equation}
  \ln x_i \;=\; z_i \;\equiv\; \frac{\mu_i^{*} - \tilde{\mu}_i^{\circ}(P,T)}{RT}.
  \label{eq:gathered_z}
\end{equation}

Imposing the constraint that the sum of the chemical end members is 1 ($\sum_{i\in\phi} x_i = 1$), it follows that

\begin{equation}
  x_i \;=\; \frac{x_i}{\sum_{j\in\phi} x_j} \;=\; \frac{e^{z_i}}{\sum_{j\in\phi} e^{z_j}}
  \;=\; \text{softmax} (\bm{z})_i .
  \label{eq:ideal_softmax}
\end{equation}

so the softmax is the exact functional form describing the equilibrium composition as a function 
of the chemical potential with respect to a reference state, normalized by $RT$. The network does not 
need to track the non-ideal terms because it predicts the logits $\hat{z}_i$ that satisfy it
directly, which absorbs the activity model into the learned mapping from
$(P, T, \text{bulk composition})$ to $\bm{z}$. In the ideal limit ($\gamma_i = 1$),
the logit output z represents a unitless chemical potential of each endmember relative to the reference state.

\subsection{A non-negative, closure-preserving reparameterization of the
pyroxene and spinel components of the \textsc{MELTS} models}
\label{sec:pxsp}

The softmax function always produces outputs that are strictly positive and sum to one, so the network
cannot emit a non-physical composition (e.g. negative atoms) for phases composed of components that are themselves 
legal compositions, regardless of its unconstrained raw output. For spinels and pyroxenes, however,
\textsc{MELTS} predicts legal compositions that are outside the simplex defined by the end members 
(i.e., some components can be negative). A softmax activation therefore cannot be applied to these components directly, 
but these compositions can be remapped to another basis that iscompatible with a softmax activation: 
components that are non-negative and sum to one.

Even for spinels and pyroxenes, which can contain negative components, there exists a limit lying outside
the simplex that expresses legal compositions, implying that there exist some compositional basis that 
can remap legal compositions to lie within the simplex. We were able to analytically derive a solution 
for pyroxene, but spinel was instead solved numerically. All other phases' compositional bases in 
\textsc{MELTS} and \textsc{HeFESTo} are already strictly non-negative.

\subsubsection{Pyroxene} We write the mole fractions of the seven \textsc{MELTS} pyroxene
components (diopside, clinoenstatite, hedenbergite, alumino-buffonite, buffonite,
essenite, jadeite) as $x_0,\dots,x_6$. Pyroxene has the general formula
M2\,M1\,T$_2$O$_6$, with one large octahedral site (M2), one small octahedral site
(M1), and two tetrahedral sites (T); the site occupancies of the seven components
follow \citet{SackGhiorso1994} and are given in Table~\ref{tab:pxsites}. The
atoms-per-formula-unit totals implied by that table are
\begin{align}
  \mathrm{Ca} &= x_0+x_2+x_3+x_4+x_5, &
  \mathrm{Na} &= x_6, \nonumber\\
  \mathrm{Mg} &= x_0+2x_1+\tfrac12(x_3+x_4), &
  \mathrm{Al} &= x_3+x_5+x_6, \label{eq:pxapfu}\\
  \mathrm{Fe}^{3+} &= x_4+x_5, &
  \mathrm{Ti} &= \tfrac12(x_3+x_4). \nonumber
\end{align}

\begin{table}[t]
\centering
\caption{Site occupancies of the seven \textsc{MELTS} pyroxene components
\citep{SackGhiorso1994}.}
\label{tab:pxsites}
\begin{tabular}{llll}
\toprule
Component & M2 & M1 & T,\,T \\
\midrule
diopside ($x_0$)            & Ca & Mg                            & Si, Si \\
clinoenstatite ($x_1$)      & Mg & Mg                            & Si, Si \\
hedenbergite ($x_2$)        & Ca & Fe$^{2+}$                     & Si, Si \\
alumino-buffonite ($x_3$)   & Ca & $\tfrac12$Mg\,$+\,\tfrac12$Ti & Al, Si \\
buffonite ($x_4$)           & Ca & $\tfrac12$Mg\,$+\,\tfrac12$Ti & Fe$^{3+}$, Si \\
essenite ($x_5$)            & Ca & Fe$^{3+}$                     & Al, Si \\
jadeite ($x_6$)             & Na & Al                            & Si, Si \\
\bottomrule
\end{tabular}
\end{table}

Seven components subject to closure leave six independent directions. These separate
naturally into three groups. Three elements are carried by exactly one component
each and can be read off directly: Fe$^{2+}$ (hedenbergite), Na (jadeite), and Mg on
the M2 site specifically (clinoenstatite, the Mg--Ca exchange vector distinct from
the M1 magnesium already carried by diopside). Three further elements---Fe$^{3+}$,
Al and Ti---are each split across two or three components, so no single component
fraction stands in for them; the corresponding coordinates are obtained by requiring
their pairwise sums to equal the element totals of Eq.~\eqref{eq:pxapfu},
\begin{equation}
  r_1+r_3 = \mathrm{Al}, \qquad r_2+r_3 = \mathrm{Fe}^{3+}, \qquad
  r_1+r_2 = 2\,\mathrm{Ti},
\end{equation}
which is a $3\times3$ system with the unique solution given below. The remaining two
quantities, Ca and Si, are never tracked directly; they are pinned by requiring each
crystallographic site to be exactly full---a second, local layer of closure nested
inside the overall constraint $\sum_i x_i = 1$. Collecting these,
\begin{equation}
  r_0=x_2,\quad r_1=x_3+\tfrac12 x_6,\quad r_2=x_4-\tfrac12 x_6,\quad
  r_3=x_5+\tfrac12 x_6,\quad r_4=x_6,\quad r_5=x_1 .
  \label{eq:pxr}
\end{equation}
Each $r_i$ is bounded for any physically legal pyroxene, but not on $[0,1]$.
Rescaling and recombining gives six quantities that are individually bounded on
$[0,1]$:
\begin{equation}
  u_0=\tfrac12 r_0,\quad u_1=\tfrac12(r_1+r_2),\quad u_2=\tfrac12(r_1+r_3),\newline
  u_3=\tfrac12(r_2+r_3),\quad u_4=\tfrac12 r_4,\quad u_5=\tfrac12 r_4+r_5 .
  \label{eq:pxu}
\end{equation}

Unit-sum closure of the full transformed vector is a separate matter and
does \emph{not} follow from site positivity. Because $\sum_i x_i = 1$, the constant 
$u_C \equiv 1 = \sum_i x_i$ is itself a linear function of the mole fractions and can 
be appended as a seventh coordinate. This makes the map square (and hence invertible), 
and gives the diopside reference endmember, for which $u_0,\dots,u_5$ all vanish, a non-zero value in the transformed space. 
The resulting vector $(u_C,u_0,\dots,u_5)$ is non-negative for any legal pyroxene but does not sum to one. 
Closure was imposed by evaluating the linear map of
Eqs.~\eqref{eq:pxr}--\eqref{eq:pxu} at each of the seven pure endmembers and dividing
each resulting row by its own row sum, yielding the fixed matrix
$\mathsf{T}_{\mathrm{px}}$ of Table~\ref{tab:pxmatrix}, in which every row is
individually non-negative and sums to unity.

\begin{table}[t]
\centering
\caption{The pyroxene reparameterization matrix $\mathsf{T}_{\mathrm{px}}$. Rows are
the seven \textsc{MELTS} pyroxene components; columns are the transformed
coordinates. Every row is non-negative and sums to unity. The same matrix is applied
to both the orthopyroxene and clinopyroxene solutions.}
\label{tab:pxmatrix}
\begin{tabular}{lccccccc}
\toprule
Component & $u_C$ & $u_0$ & $u_1$ & $u_2$ & $u_3$ & $u_4$ & $u_5$ \\
\midrule
diopside          & 1   & 0   & 0   & 0   & 0   & 0   & 0 \\
clinoenstatite    & 1/2 & 0   & 0   & 0   & 0   & 0   & 1/2 \\
hedenbergite      & 2/3 & 1/3 & 0   & 0   & 0   & 0   & 0 \\
alumino-buffonite & 1/2 & 0   & 1/4 & 1/4 & 0   & 0   & 0 \\
buffonite         & 1/2 & 0   & 1/4 & 0   & 1/4 & 0   & 0 \\
essenite          & 1/2 & 0   & 0   & 1/4 & 1/4 & 0   & 0 \\
jadeite           & 2/5 & 0   & 0   & 1/5 & 0   & 1/5 & 1/5 \\
\bottomrule
\end{tabular}
\end{table}

Because closure was imposed by endmember normalization rather than derived from the
site-occupancy argument, its validity for a general (non-endmember) composition is
not guaranteed \emph{a priori}, and was instead established empirically: applying
$\mathsf{T}_{\mathrm{px}}$ to the full set of \textsc{MELTS}-computed pyroxene
compositions in the training data yields vectors that are non-negative and sum to
unity, and $\mathsf{T}_{\mathrm{px}}^{-1}$ recovers the original mole fractions 
within numeric precision.

\subsubsection{Spinel} We were unable to reduce the spinel components (chromite, hercynite, magnetite, spinel
\emph{sensu stricto}, ulv\"ospinel) to a comparably clean set of site-occupancy inequalities. No analogous
analytic transform was therefore pursued. Instead a masked, invertible linear
transform $\mathsf{T}_{\mathrm{sp}}$ was fit numerically.

The matrix was parameterized as lower triangular, which guarantees invertibility for
any non-zero diagonal. Chromite uniquely expresses Cr contents and is therefore always 
non-negative. It was mapped onto its own dedicated output coordinate and excluded from the fit
so that chromium-free models would be able to use the same transform. The
remaining entries were optimized by gradient descent (Adam; PyTorch) against
approximately $3\times10^{6}$ \textsc{MELTS} spinel compositions drawn from the
training data, minimizing a two-term objective: a mean-squared term driving each
transformed composition's row sum toward unity, and a leaky-ReLU penalty on negative
transformed entries driving them toward zero. The learning rate was annealed by a
factor of ten every 2{,}500 of 10{,}000 steps; the diagonal was reset to unity
whenever it decayed toward zero, preserving invertibility; and coefficients falling
below $0.05$ in magnitude were periodically hard-thresholded to zero to keep the
fitted matrix sparse. The solution satisfyingly converged to one that can be expressed 
in terms of relatively simple fractions.

\begin{table}[t]
\centering
\caption{The spinel reparameterization matrix $\mathsf{T}_{\mathrm{sp}}$. Rows are
the four (non-chromite) \textsc{MELTS} spinel components; columns are the transformed
coordinates. Every row is non-negative and sums to unity.}
\label{tab:spmatrix}
\begin{tabular}{lcccc}
\toprule
Component & $u_C$ & $u_0$ & $u_1$ & $u_2$ \\
\midrule
Hercynite     & 1   & 0   & 0   & 0 \\
Magnetite    & 2/5 & 3/5   & 0   & 0 \\
Spinel      & 19/20 & 0  & 1/20   & 0 \\
Ulv\"ospinel & 1/5 & 0   & 0  & 4/5 \\
\bottomrule
\end{tabular}
\end{table}

The resulting matrix and its analytic inverse were empirically validated in the same way as the
pyroxene transform: by histogramming the transformed compositions and their row
sums, checking their extrema, and comparing round-tripped compositions against the
originals across the full training set. Original compositions are recovered within numeric precision 
and the transformed components of all \textsc{MELTS} spinels are confined almost entirely to $[0,1]$.

============================================================
\subsection{Legality polish for pyroxene and spinel compositions}
\label{sec:polish}
%=====================================================================================

The transforms derived above in \S\ref{sec:pxsp} guarantees that every legal composition lands in the simplex 
($\hat x \ge 0$, $\sum\hat x = 1$), but not the reverse: the legal compositions fill only 
part of the simplex, and other regions of the simplex are forbidden for legal compositions. 
Because the softmax can return any point in the simplex, it is possible 
for a softmax output to be illegal without further constraints. To address this, 
we apply a post-processing step to the softmax outputs of pyroxene and spinel 
to force them to be legal compositions. This correction is commonly applied to orthopyroxene (6-7\% of predicted OPX), for which
CaO is low generally and can easily be pushed negative. For spinel and clinopyroxene, these corrections are more rare, affecting less than 2\% 
of predicted spinel and 0.1\% of predicted clinopyroxene. 


\begin{table}[!htb]
\centering
\footnotesize
\setlength{\tabcolsep}{3pt}
\renewcommand{\arraystretch}{1.15}
\begin{tabular}{@{}l l r r r r r r r r r@{}}
\toprule
Component & Definition ($T^{-1}$) & Al$_2$O$_3$ & CaO & FeO & Fe$_2$O$_3$ & MgO & Na$_2$O & SiO$_2$ & TiO$_2$ & Cr$_2$O$_3$ \\
\midrule
\multicolumn{11}{@{}l}{\emph{Pyroxene}}\\
$u_C$ & Di                          & 0   & 1  & 0 & 0 & 1  & 0   & 2 & 0 & -\\
$u_0$ & $3$Hd $-$ $2$Di             & 0   & 1  & 3 & 0 & $-2$ & 0 & 2 & 0 & -\\
$u_1$ & $2$Abf $+$ $2$Bf $-$ $2$Ess $-$ Di & 0 & 1 & 0 & 0 & 1 & 0 & 0 & 2 & -\\
$u_2$ & $2$Abf $-$ $2$Bf $+$ $2$Ess $-$ Di & 2 & 1 & 0 & 0 & $-1$ & 0 & 0 & 0 & -\\
$u_3$ & $2$Bf $+$ $2$Ess $-$ $2$Abf $-$ Di & 0 & 1 & 0 & 2 & $-1$ & 0 & 0 & 0 & -\\
$u_4$ & $5$Jd $+$ $2$Bf $-$ $2$En $-$ $2$Abf $-$ $2$Ess & 0.5 & $-2$ & 0 & 0 & $-4$ & 2.5 & 4 & 0 & -\\
$u_5$ & $2$En $-$ Di                & 0   & $-1$ & 0 & 0 & 3 & 0 & 2 & 0 & -\\
\midrule
\multicolumn{11}{@{}l}{\emph{Spinel}}\\
Chr     & Chr                                        & 0   & -  & 1     & 0   & 0  & -    & - & 0 & 1 \\
$u_0$   & Hc                                         & 1   & -  & 1     & 0   & 0  & -    & - & 0 & 0\\
$u_1$   & $\tfrac53$Mt $-$ $\tfrac23$Hc              & $-\tfrac23$ & -  & 1 & $\tfrac53$ & 0 & - & - & 0 & 0\\
$u_2$   & $20$Sp $-$ $19$Hc                          & 1   & - & $-19$ & 0   & 20 & -    & - & 0 & 0\\
$u_3$   & $\tfrac54$Usp $-$ $\tfrac14$Hc           & $-\tfrac14$ & - & $\tfrac94$ & 0 & 0 & - & - & $\tfrac54$ & 0\\
\bottomrule
\end{tabular}
\caption{Transformed components and their chemical definitions. These are the rows of the 
inverted transform matrices $\mathsf{T}_{\mathsf{px}}^{-1}$ and $\mathsf{T}_{\mathsf{sp}}^{-1}$.
Mapped into an oxide space, these transformed component contain negative atoms and therefore a 
softmax output in this space is capable of producing illegal negative atoms. 
Endmember abbreviations: Di diopside, En clinoenstatite,
Hd hedenbergite, Abf alumino-buffonite, Bf buffonite, Ess essenite, Jd jadeite;
Chr chromite, Hc hercynite, Mt magnetite, Sp spinel (MgAl$_2$O$_4$), Usp ulv\"ospinel.}
\label{tab:Kcomp}
\end{table}

\subsubsection{Correcting a single negative oxide}

A single negative oxide can be safely corrected with trivial arithmetic. 

The first step in correcting a single negative oxide is to separate the positive and negative oxide contributions from each component ($S_+$ and $S_-$ respectively). 
For oxide $i$, let $q_i$ be the column of oxide stoichiometry for the transformed components as given in Table~\ref{tab:Kcomp}. 
\begin{equation}
  S_+ = \sum_{q_i > 0} q_i\,\hat x_i, \qquad
  S_- = \sum_{q_i < 0} |q_i|\,\hat x_i, \qquad
  \text{oxide moles} = S_+ - S_- 
\end{equation}
The oxide is negative if $S_- > S_+$. We employ a simple fix that aims to preserve as much of the original prediction as possible. 
We find two scaling factors to scale the positive and negative contributing components to balance the oxide. We multiply every component fraction on
the positive side by one factor $a$ and every fraction on the negative side by one factor $b$. If we set the 
oxide to zero and enforce the closure constraint, we formulate a system of two equations that 
fix the two unknowns $a$ and $b$:
\begin{align}
  \text{oxide becomes exactly zero:}\quad & a S_+ = b S_-, \label{eq:zero}\\
  \text{fractions still sum to one:}\quad & a P + b N = 1, \label{eq:close}
\end{align}
where $P$ and $N$ are the plain sums of the component fractions on each side.
\begin{equation}
  P = \sum_{\hat x_i > 0} \hat x_i, \qquad
  N = \sum_{\hat x_i < 0} \hat x_i, \qquad
  P + N = 1
\end{equation}

From Eq.~\eqref{eq:zero}, $b = a S_+/S_-$; substituting into
Eq.~\eqref{eq:close} gives
\begin{equation}
  a = \frac{S_-}{S_- P + S_+ N}, \qquad
  b = \frac{S_+}{S_- P + S_+ N}.
  \label{eq:ab}
\end{equation}

\subsubsection{Pyroxene: CaO $< 0$}

From Table~\ref{tab:Kcomp},
\begin{equation}
  S_+ = \hat x_C + \hat x_0 + \hat x_1 + \hat x_2 + \hat x_3, \qquad
  S_- = 2\hat x_4 + \hat x_5 .
\end{equation}
All positive coefficients are 1, so $P = S_+$ and $N = \hat x_4 + \hat x_5$, and
Eq.~\eqref{eq:ab} simplifies to:
\begin{equation}
  b = \frac{1}{3\hat x_4 + 2\hat x_5}, \qquad
  a = \frac{S_-}{S_+\,(3\hat x_4 + 2\hat x_5)} .
\end{equation}

The CaO correction is applied to clinopyroxene only.

\subsubsection{Pyroxene: MgO $<0$}

\begin{equation}
  S_+ = \hat x_C + \hat x_1 + 3\hat x_5, \quad P = \hat x_C + \hat x_1 + \hat x_5; \qquad
  S_- = 2\hat x_0 + \hat x_2 + \hat x_3 + 4\hat x_4, \quad N = \hat x_0 + \hat x_2 + \hat x_3 + \hat x_4,
\end{equation}
$a$, $b$ then can be calculated using Eq.~\eqref{eq:ab}. The MgO correction
is applied to clinopyroxene only.

\subsubsection{Spinel: FeO $< 0$ and Al$_2$O$_3$ $< 0$}

Spinel has two oxides that can go negative, sometimes simultaneously, substantially complicating the problem. From Table~\ref{tab:Kcomp},
\begin{align}
  \text{FeO} &= \hat x_{\mathsf{Chr}} + \hat x_0 + \hat x_1 + \tfrac94\hat x_3 - 19\,\hat x_2, \label{eq:spFe}\\
  \text{Al}_2\text{O}_3 &= \hat x_0 + \hat x_2 - \tfrac23\hat x_1 - \tfrac14\hat x_3 . \label{eq:spAl}
\end{align}
In models without Cr, $\hat x_{\mathsf{Chr}} = 0$.

\paragraph{Three pools.} Because the two constraints share components, the components are
grouped by the signs they carry in the two equations, and each group gets its own factor:
\begin{center}
\begin{tabular}{@{}l l c c c@{}}
\toprule
Pool & Components & FeO & Al$_2$O$_3$ & Factor \\
\midrule
$A = \hat x_{\mathsf{Chr}} + \hat x_0$ & Chr, $u_0$ & $+$ & $+$ & $a$ \\
$B = \hat x_1 + \hat x_3$         & $u_1$, $u_3$ & $+$ & $-$ & $b$ \\
$C = \hat x_2$                    & $u_2$ & $-$ & $+$ & $c$ \\
\bottomrule
\end{tabular}
\end{center}
In all cases, the closure constraint is enforced:
\begin{equation}
  aA + bB + cC = 1, \label{eq:spclose}
\end{equation}
If only one oxide is negative, then only the violated oxide is set to zero, 
and the pool it does not involve held fixed.
Then, the second equation for each scenario is:
\begin{itemize}[nosep]
  \item \emph{FeO only} ($b = 1$): Eq.~\eqref{eq:spclose} and
        $aA + \hat x_1 + \tfrac94\hat x_3 = 19\,cC$;
  \item \emph{Al$_2$O$_3$ only} ($c = 1$): Eq.~\eqref{eq:spclose} and
        $a\hat x_0 + C = b\big(\tfrac23\hat x_1 + \tfrac14\hat x_3\big)$.
\end{itemize}

\paragraph{Both violated.} If both oxides are negative, or if fixing one would push the other
below zero, all three factors are solved together with three constraints: closure, FeO $= 0$, and
Al$_2$O$_3$ $= 0$,
\begin{equation}
  \begin{pmatrix}
    A & B & C\\
    A & \hat x_1 + \tfrac94\hat x_3 & -19C\\
    \hat x_0 & -\big(\tfrac23\hat x_1 + \tfrac14\hat x_3\big) & C
  \end{pmatrix}
  \begin{pmatrix} a\\ b\\ c \end{pmatrix}
  =
  \begin{pmatrix} 1\\ 0\\ 0 \end{pmatrix}.
\end{equation}

\todo{Add Spinel 1:1 figure}

%=====================================================================================
\subsection{Mass balance}
\label{sec:massbalance}
%=====================================================================================

The assemblage predicted by the emulator is not strictly constrained to reproduce the bulk composition;
it is only encouraged to with a loss term during training (\S\ref{ch4-training}). 
Therefore, a post-processing step is necessary to ensure that the final assemblage is consistent with the requested bulk composition.
\emph{nGibbs} uses a mass balance correction adapted from \citet{asimow1998}, who find the smallest relative 
change to $\hat m$ that makes $\hat b = b$. The implementation in \emph{nGibbs} differs from their procedure most
strongly in two ways: (1) the minimisation is solved in parallel for thousands of assemblages at once on a GPU, and (2)
the correction is done iteratively, with constraints enforced at each step so that no elemental abundance becomes negative. 

The section makes use of variables defined in Table~\ref{tab:key} as well as those unique to this section
in Table~\ref{tab:mbkey}. Index forms are given alongside the matrix forms
throughout, with $c$ running over components and $e, e'$ over elements.

\begin{table}[!htb]
\centering
\small
\renewcommand{\arraystretch}{1.25}
\begin{tabular}{@{}c c >{\raggedright\arraybackslash}p{11.6cm}@{}}
\toprule
Symbol & Dim. & Meaning \\
\midrule
$\delta b$ & E & Bulk residual in element moles, $\delta b = \Gamma\, b - \hat m A$ \\
$\delta \hat m$ & C & Correction to the component moles \\
$\zeta$ & C & Relative correction, $\zeta_c = \delta\hat m_c / \hat m_c$ \\
$D$ & C$\times$E & Abundance-weighted stoichiometry, $D = \operatorname{diag}(\hat m)\, A$,
  i.e.\ $D_{ce} = \hat m_c A_{ce}$ \\
$\nu$ & E & Lagrange multipliers, one per element \\
$G$ & E$\times$E & $G = D^{\!\top} D$, the matrix of the system solved for $\nu$ \\
$\kappa$ & 1 & Regularisation added to the diagonal of $G$ ($10^{-12}$) \\
$\sigma_i,\ \Phi,\ \Psi$ & E, C$\times$E, E$\times$E & Singular values and singular vectors
  of $D$: $D = \Phi\, \operatorname{diag}(\sigma)\, \Psi^{\!\top}$ \\
$N_{\rm it}$ & 1 & Number of correction iterations (3) \\
\bottomrule
\end{tabular}
\caption{Symbols added for the mass balance. Dimensions as in Table~\ref{tab:key}.}
\label{tab:mbkey}
\end{table}

The neural network returns the component moles $\hat m$ (the endmembers of every phase, in the
transformed basis of Section~\ref{sec:polish}). The element moles they imply are $\hat m A$
($[1\times\text{C}][\text{C}\times\text{E}] \to [1\times\text{E}]$). The reconstructed bulk is normalized to sum to one, then
it's residual is given by
\begin{equation}
  \delta b = b - \hat m A 
  \qquad .
  \label{eq:mbr}
\end{equation}

Mass balance seeks a correction $\delta\hat m$ to the component moles such that
\begin{equation}
  \delta\hat m\, A = \delta b,
  \qquad\text{i.e.}\qquad
  \sum_c \delta\hat m_c\, A_{ce} = \delta b_e \quad \text{for each element } e .
  \label{eq:mbcond}
\end{equation}

This is E equations (about ten) in C unknowns (several tens), so the system is underconstrained and 
has infinitely many solutions. Because the network's prediction is the best available estimate of the assemblage, it is 
most reasonable to choose the solution that changes it least. We define this as the solution that minimizes the sum of 
squared relative changes to the component moles. For convenience, we define the relative change $\zeta_c$ for each component $c$ as

\begin{equation}
  \zeta_c = \frac{\delta\hat m_c}{\hat m_c},
  \label{eq:zeta}
\end{equation}
We minimize the sum of squared relative changes,
\begin{equation}
  \tfrac12\, \zeta\, \zeta^{\!\top}
  = \tfrac12 \sum_c \left(\frac{\delta\hat m_c}{\hat m_c}\right)^{\!2}.
  \label{eq:mbobj}
\end{equation}
The factor $\tfrac12$ is arbitrarily added to elegantly cancel the 2 produced by the differentiation below. We introduce
a new variable: $D$, the component abundance-weighted stoichiometry matrix, to relate $\zeta$ to the bulk residual $\delta b$.
\begin{equation}
  \zeta\, D = \delta b,
  \qquad
  D = \operatorname{diag}(\hat m)\, A \quad (\text{C}\times\text{E}),
  \label{eq:agweighted}
\end{equation}
A component absent from the assemblage ($\hat m_c = 0$) has a zero row at $D_c$ and, by Eq.~\eqref{eq:agback}, $\delta\hat m_c = 0$ 
so absent components are never changed.

We therefore formulate the problem as a minimization of Eq.~\eqref{eq:mbobj} subject to the constraints of Eq.~\eqref{eq:agweighted}. 
In other words: our solution is the smallest relative change to the component moles necessary to satisfy conservation of mass. 

\subsubsection{Solving with Lagrange multipliers}
\label{sec:mb-lagrange}

Lagrange multipliers reformulate this problem into a set of ordinary equations. One new unknown, the
multiplier $\nu_e$, is introduced for each element. Each constraint is written as an expression that must equal zero.
The multipliers $\nu_e$ relate the derivatives of the objective function to the constraints. Requiring every derivative 
of the result to vanish then expresses two conditions at once: the constraints hold, and no change that respects them
can lower the objective $ \tfrac12\, \zeta\, \zeta^{\!\top}$. This is the Lagrangian,
\begin{equation}
  \mathcal{L}(\zeta, \nu)
  = \tfrac12\, \zeta\,\zeta^{\!\top} - \big( \zeta D - \delta b \big)\, \nu^{\!\top},
  \label{eq:lagrangian}
\end{equation}
with $\nu$ the row of E multipliers. Now, differentiating with respect to the
relative correction $\zeta$ gives
\begin{equation}
  \frac{\partial \mathcal{L}}{\partial \zeta} = \zeta - \nu\, D^{\!\top} = 0
  \qquad\Longrightarrow\qquad
  \zeta = \nu\, D^{\!\top}
  \quad [1\times\text{E}][\text{E}\times\text{C}] \to [1\times\text{C}].
  \label{eq:ystat}
\end{equation}
This reveals that the C relative corrections are not independent: they are all determined
by just E numbers, the Lagrange multipliers. The problem has been reduced from C unknowns to E,
and now the system is well-constrained.

Substituting Eq.~\eqref{eq:ystat} into the $\zeta\, D = \delta b$ constraint gives
\begin{equation}
  \nu\, D^{\!\top} D = \delta b,
\end{equation}

and now the only unknowns are the E multipliers $\nu$, which can solved with an LU (lower-upper) decomposition as 

\begin{equation}
  \nu^{\!\top} = SOLVE(D^{\!\top} D, \delta b^{\!\top}) .
  \label{eq:nu}
\end{equation}

With $\nu$ known, Eq.~\eqref{eq:ystat} relates the relative correction $\zeta$ to the multipliers.

\begin{equation}
  \zeta = \nu A^{\!\top} \operatorname{diag}(\hat m)
\end{equation}

The absolute correction $\delta\hat m$ is then obtained by multiplying the relative correction by the component moles (Eq.~\eqref{eq:zeta}).

\begin{equation}
  \delta\hat m = \zeta \operatorname{diag}(\hat m) = \nu A^{\!\top} \operatorname{diag}(\hat m)^2
  \quad [1\times\text{E}][\text{E}\times\text{C}][\text{C}\times\text{C}] \to [1\times\text{C}],
  \label{eq:share}
\end{equation}

\paragraph{Why the Lagrange form is used.} \citet{asimow1998} obtain the same solution by
SVD, however their problem is distinct from this one in that they correct one assemblage at a
time, whereas \emph{nGibbs} applies this operation to tens of thousands of assemblages in parallel. 
The Lagrange form is used here because batched SVD handles these matrices poorly on a GPU.
An informal test on a dataset of 15,000 assemblages from the HeFESTo Earth model validation dataset (E = 8, C = 68) 
shows HeFESTo test bundle (B = 15{,}000, E = 8, C = 68), \code{torch.linalg.pinv} took 10.1\,s on the GPU against 0.29\,s 
on the CPU (single precision), whereas the Lagrange solve took 25\,ms on the GPU (double precision; 48\,ms on
the CPU). The Lagrange form, however, requires double precision floating point arithmetic to recover an accurate solution,
whereas the SVD is robust to single precision arithmetic because the $\nu\left( D^{\!\top} D \right)$ solve of the Lagrange formulation
squares the condition number. This cost is marginal, however, beside the speedup of the Lagrange form on both CPU and GPU, and it elegantly
avoids the need to bus data between CPU and GPU when \emph{nGibbs} is run on a GPU.
\subsection{Regularisation}
\label{sec:mb-reg}

\paragraph{The regularized solve.} In practice, a small number $\kappa$ is added to every diagonal entry
of $\left( D^{\!\top} D \right)$,
\begin{equation}
  \nu\,\big( D^{\!\top} D  + \kappa I\big) = \delta b,
  \qquad \kappa = 10^{-12},
  \label{eq:regsolve}
\end{equation}
where $I$ is the E$\times$E identity matrix. This is because $ D^{\!\top} D $ has all-zero columns and rows for 
any bulk composition lacking any particular element and is therefore singular. For a row lacking an element, 
its equation becomes $\kappa\,\nu_e = 0$ for the regularized case and $\nu$ is zero because $\delta b_e = 0$ for an absent element. 
Without regularization, the equation for this row would be $0 \nu_e = 0$, for which $\nu_e$ is undefined. 

On the \textsc{HeFESTo} Earth Model validation dataset, the smallest non-zero diagonal entries of $ D^{\!\top} D $ are $ \thicksim 3\times10^{-7}$ 
(median over rows), so $\kappa$ must be far smaller than that. A $\kappa$ as small as 10\textsuperscript{-12} requires
double precision floating point values, but this is also required by the solve of the Lagrange multipliers and therefore poses no additional complication.

% Generated by Chapter4/Figures/fig_mass_balance_G_hist.py
\begin{figure}[htbp]
  \centering
  \includegraphics[width=8.3cm]{mass_balance_G_hist_DtD_diag_linear.pdf}
  \caption{Histogram of the non-zero diagonal entries of $D^{\!\top} D$ formed by the
  conservation of mass correction (all three iterations) on the 15{,}000 assemblages of the
  HeFESTo Earth adiabat isentropic (NPS) test set. The dashed line marks $\kappa = 10^{-12}$.}
  \label{fig:mb-dtd-diag}
\end{figure}

\subsection{Seeding phases for unexplained elements}
\label{sec:seed}

\paragraph{Why it is needed.} The conservation of mass correction can only move an element through components
that contain it and are present ($\hat m_c > 0$). If the predicted assemblage contains no
carrier of an element that is in the bulk (for example K, P or H when the network places
the K-, P- or H-bearing phase just below saturation), that element's residual cannot be
removed: the regularised solve leaves it in place (Section~\ref{sec:mb-reg}). In the 102
Closed Cr test set this happens in 13 of 15{,}000 rows.

\paragraph{Procedure.} Before the correction, for the continuous models
(\code{NN_MELTS.seed_unexplained_phases}), repeat for each row:
\begin{enumerate}[nosep]
  \item \textbf{Unexplained elements.} Element $e$ is unexplained if $b_e > 0$ and no
        component with $\hat m_c > 0$ contains it ($A_{ce} = 0$ for every present $c$). Call
        the set of unexplained elements $\mathcal{U}$.
  \item \textbf{Candidates.} Absent phases ($\hat n_p = 0$) whose \emph{predicted}
        composition contains at least one element of $\mathcal{U}$. The chemistry heads
        predict a composition $\hat x$ for every phase, present or not, so one is always
        available.
  \item \textbf{Choice.} The candidate closest to saturation: the largest (least negative)
        affinity $g_p$ ($\hat g_p$ for \code{AffinityModel}).
  \item \textbf{Abundance.} Let
        \begin{equation}
          a_e = \sum_c M_{pc}\,(\hat x V + f)_c\, A_{ce}
          \label{eq:seedcomp}
        \end{equation}
        be the moles of element $e$ in one mole of the chosen phase $p$: its predicted
        component fractions (Eq.~\ref{eq:m}, restricted to phase $p$ by $M_{pc}$) times
        their stoichiometry. The phase is added with the abundance $\hat n_p$ that best
        supplies the elements of $\mathcal{U}$ it carries,
        \begin{equation}
          \hat n_p = \frac{\sum_{e \in \mathcal{U}} a_e\, \Gamma\, b_e}{\sum_{e \in \mathcal{U}} a_e^2},
          \label{eq:seed}
        \end{equation}
        the least-squares fit of $\hat n_p\, a_e = \Gamma\, b_e$ over those elements. When the
        phase covers a single unexplained element, $\hat n_p$ supplies exactly the bulk's
        amount of it, $\hat n_p = \Gamma\, b_e / a_e$.
  \item \textbf{Insert.} The phase's components get $\hat m_c = \hat n_p\,(\hat x V + f)_c$
        for every $c$ with $M_{pc} = 1$, as in Eq.~\eqref{eq:m}.
\end{enumerate}
The loop stops when $\mathcal{U}$ is empty or no candidate is left; each round adds a new
phase, so it ends after at most P rounds. $\hat m$ is then renormalised so that
$\Gamma = 1$. Rows that need no seeding are not touched.

\paragraph{Effect.} Seeding gives every element of the bulk a non-zero column in $D$.
Because the seeded phase already carries about the right amount of the missing element, it
also gets a sensible weight: the correction moves each component in proportion to
$\hat m_c^2$ (Eq.~\ref{eq:share}), so a phase added at a negligible abundance would barely
be used. On the 102 test set, 14 phases were seeded in the 13 rows, and 11 of them are
present in the ground truth (apatite, leucite, fluid, hornblende). A single unclamped
correction step on those rows reduces their residuals from order $10^{-3}$ to about
$10^{-8}$.

\subsection{The iterative correction}
\label{sec:mb-iter}

\paragraph{Why one step is not enough.} The step of Sections~\ref{sec:mb-lagrange}--\ref{sec:mb-reg}
places no bound on $\delta\hat m_c$. Nothing prevents $\delta\hat m_c < -\hat m_c$, which
would leave a component with negative moles; this is the ``overstepping'' failure Asimow
and Ghiorso note. Setting such a component to zero restores a physical assemblage but
breaks the mass balance again, because the moles it would have supplied are now missing.
The step is therefore repeated: each repetition spreads the remaining residual over the
components that are still present.

\paragraph{Procedure.} Starting from the seeded $\hat m$, each of $N_{\rm it}$ iterations
($N_{\rm it} = 3$ by default) does:
\begin{enumerate}[nosep]
  \item \textbf{Residual on the current scale.} Recompute $\Gamma = \sum_e (\hat m A)_e$ and
        $\delta b = \Gamma\, b - \hat m A$ (Eq.~\ref{eq:mbr}). Only the proportions of the
        bulk are enforced, and each clamp (step 4) changes the total, so the target
        $\Gamma\, b$ is rebuilt every iteration. A target held at a fixed scale would make
        the residual grow between iterations.
  \item \textbf{Weights.} Build $D = \operatorname{diag}(\hat m)\,A$, with each present
        $\hat m_c$ raised to at least $10^{-9}$ in the diagonal. The floor keeps a very small
        but present component solvable; absent components keep a zero row.
  \item \textbf{Solve.} Form $G = D^{\!\top} D$, solve the regularised system
        Eq.~\eqref{eq:regsolve} for $\nu$, then $\zeta = \nu D^{\!\top}$
        (Eq.~\ref{eq:ystat}), all in double precision. This is the operation that runs
        batched on the GPU.
  \item \textbf{Update and clamp.} $\hat m \leftarrow \max\big(\hat m + \zeta \operatorname{diag}(\hat m),\ 0\big)$,
        element by element: Eq.~\eqref{eq:agback}, followed by setting any negative
        component moles to zero.
  \item \textbf{Polish.} For MELTS models containing orthopyroxene, clinopyroxene or spinel,
        re-impose the legality corrections of Section~\ref{sec:polish}. The update moves
        component moles within a phase and can create a negative oxide in a phase whose
        predicted composition was legal.
\end{enumerate}
After the last iteration the remaining residual, $\lVert \hat b - b \rVert$ with $\hat b$ from
Eq.~\eqref{eq:bhat}, is reported for each row.

\paragraph{Properties.}
\begin{itemize}[nosep]
  \item \emph{Termination.} A component set to zero has a zero row in $D$ from then on and
        cannot return. The set of present components can only shrink, so the iteration
        cannot cycle, and a phase removed by the correction stays removed.
  \item \emph{Exactness.} Each unclamped step removes the residual up to the
        regularisation. After the final clamp and polish a small residual can remain: the
        median on the 102 test set is $3.0\times10^{-8}$.
  \item \emph{Non-negativity and legality.} Returned component moles are never negative,
        and the polish runs after every step rather than only at the end.
  \item \emph{Special cases.} With one iteration, no clamp, no polish and $\kappa \to 0$,
        the procedure reduces to the single weighted step of Asimow and Ghiorso,
        Eq.~\eqref{eq:agsvd}. That single step remains available for comparison as
        \code{mass_balance='pinv'}, evaluated by SVD on the CPU. Setting
        \code{relative=False} replaces $\operatorname{diag}(\hat m)$ by the identity on
        present components, giving the unweighted smallest change.
\end{itemize}


\subsection{Data generation}
\label{sec:datagen}

\subsubsection{Bulk composition sampling}

Bulk compositions are sampled from whole-rock analyses in the GEOROC and PetDB
databases \citep{GEOROC,PetDB}, preferentially weighting the compositional range
the emulator is intended to serve. Whole rock analyses from GEOROC in particular are weighted towards more evolved compositions, so mafic and ultramafic compostions were more heavily weighted in the  \emph{nGibbs} training datasets than in the original source datasets.  For the \textsc{HeFESTo} models, one fifth of the training data are drawn from the population of compositions between 5 and 20 wt\% MgO, three tenths of the data are drawn from the population of compositions with greater than 20 wt\% MgO, and the remaining half is a mixture of these two mafic and ultramafic populations in random proportions. These base compositions were substantially randomly altered so that \emph{nGibbs} would perform well in a wider variety of parameter spaces, including rock compositions that are uncommon or unknown on Earth's surface. For the Earth model, one quarter of the compositions are not perturbed. For the \textsc{MELTS} emulators, half of the training data is generated from any composition in the compiled database, one fifth of the data are pulled from ultramafic compositions with greater than 25 wt\% MgO, and three tenths of the data are pulled from compositions with greater than 5  wt\% MgO. These bulk compositions are perturbed by guassian noise with a standard deviation of 5\%, a substantially more modest perturbation that that applied to the \textsc{HeFESTo} emulator training data.  

For volatile-focused \textsc{MELTS} 1.2.0 emulators in particular, 35\% of the training and validation datasets are drawn from sedimentary compositions within the Gard database \citep{Gard2019}. Of these sedimentary bulk compositions, about one quarter are drawn from carbonate rich compositions, with CO\textsubscript{2} contents greater than 10 wt\%, whereas the remaining sedimentary compositions contain less than 10 wt\% CO\textsubscript{2}. The remaining 65\% of the training data were drawn from
igneous whole rock data from GEOROC and PetDB, similarly to the other models \citep{GEOROC,PetDB}. 

In order to train emulators that perform well in a variety of redox conditions, 
each sampled composition is randomly assigned an oxygen fugacity from -5 to +5 log$_{10}\Delta$FMQ for \textsc{MELTS} training data, or -7 to +3 log$_{10}\Delta$FMQ for \textsc{HeFESTo} training data. Open 
system models simply carry this oxygen fugacity as a feature, whereas closed models translate this oxygen fugacity into a molar \textsuperscript{+3}Fe/\textsuperscript{+2}Fe ratio using an empirical parameterization. This is not meant to be explicitly correct; it is meant to generate training data with an appropriate distribution of molar \textsuperscript{+3}Fe/\textsuperscript{+2}Fe ratios.

\begin{equation}
^{+3}Fe/^{+2}Fe \approx 10^{\left(0.2 \times \text{log}_{10}\Delta\text{FMQ} \right) - 1}
\label{eq:fO2ratio}
\end{equation}

Manganese and nickel are excluded from the compositional basis of all released
models. Manganese and nickel are not handled by \textsc{HeFESTo}, and we have 
observed that \textsc{MELTS} solutions that carry MnO and NiO frequently stabilize
phases with compositions rarely observed in nature, and there is little value in
emulating the regions of a model that are not themselves trusted. Excluding these
oxides also removes two dimensions from the component basis at small cost to
the applications the emulators target.

\subsubsection{State-variable sampling}\label{stae-var-sampling}

For the \textsc{MELTS} 1.0.2 and 1.2.0 emulator families, each sampled composition is isobarically cooled from 1800\textdegree C in 1\textdegree C increments at a random pressure between 1 and 12000 bars. These computed assemblages have well-defined entropies and therefore are acceptable for use in training isentropic models. 

\textsc{HeFESTo} Earth and Mars models are each trained around a core of (approximately) adiabatic simulations at reasonable potential temperatures (1200--1650\textdegree C for Earth; 500--1750\textdegree C for Mars). These adiabats are supplemented with cooler conditions due to a conductive lid at the surface, which is assumed to be between 0-250 km thick for Earth and 100-700 km thick for Mars.  \todo{Potential temperature citations?}. Similarly, warmer conditions are modeled at the approach to the core-mantle boundary (CMB) for the Earth model, which assume a mantle temperature between 3200--4700\textdegree C at the CMB and a 100-300km width conductive layer

Although \textsc{HeFESTo} natively calculates isentropes, such a calculation requires an iteration of T at every step to retrieve the imposed entropy, which roughly doubles the computation time. Therefore, adiabatic PT paths were instead approximated to reduce computational expense in the generation of the training dataset. To accomplish this, a second-order linear regression of temperature as a function of specific entropy and pressure was fit to a dataset of dozens of \textsc{HeFESTo} isentropes at representative compositions and potential temperatures. A similar linear regression was fit for specific entropy as a function of temperature at P = 0 GPa on the same dataset to translate potential temperatures to specific entropy. Similarly to the iron speciation parameterization described in \S\ref{state-var-sampling}, this approach is not intended to be explicitly correct; it is intended to control the distribution of the training dataset in pressure-temperature-entropy space without unnecessary computational expense. These regressions are used to impose a roughly adiabatic pressure-temperature path for training data. These paths are pertubed with small random noise to prevent the emulators from learning relationships between composition and entropy. 



\subsubsection{Execution}

\textsc{MELTS} simulations are driven through \textsc{alphaMELTS} on a window
\citep{Smith2005} run with GNU parallel. \textsc{HeFESTo} simulations are
run on the Caltech High Performance Computing cluster as SLURM array jobs. 
\todo{describe hefesto resampling procedure}

% -- PLACEHOLDER: campaign scale ---------------------------------------------
% Number of simulations and retained rows per released model.
% Include MELTStable rows
% For HeFESTo: Conductive vs adiabatic sims?
% Attrition rates
% Filtered rates
% ----------------------------------------------------------------------------

\subsubsection{Import and attrition}

Reference-code output is parsed into a single wide table in which each row is one
equilibrium state and columns follow a \texttt{quantity(phase)} convention. Retained
per row are the state variables, the bulk composition, the phase assemblage, the
per-phase component molar abundances, the per-phase masses, and the bulk physical
properties reported by the reference code. The last of these are retained for
validation only; consistent with Section~\ref{sec:overview}, they are never training
targets.

A fraction of simulations fail to converge or terminate without usable output. These
are detected and discarded at import. Attrition differs sharply between codes and
compositional domains, and is reported here because it determines the difference
between the sampled and the achieved distribution: approximately $0.1\%$ for
\textsc{HeFESTo}, approximately $15\%$ for \textsc{MELTS} on igneous bulk
compositions, and approximately $40\%$ for \textsc{MELTS} 1.2 on sedimentary bulk
compositions. The last figure reflects that \textsc{MELTS} 1.2 was applied near the
edge of its calibration.

\subsubsection{Coverage}

The achieved envelope of each released model in pressure, temperature or entropy,
and composition defines its domain of applicability.

% -- FIGURE PLACEHOLDERS -----------------------------------------------------
% \begin{figure}[t] ... \caption{Training-set density in P--T (or P--S) for each
% released model.} \label{fig:coverage-pt} \end{figure}
%
% \begin{figure}[t] ... \caption{Bulk-composition coverage of the training sets,
% projected onto Si--Mg, with the sampled database compositions for reference.}
% \label{fig:coverage-comp} \end{figure}
% ----------------------------------------------------------------------------

The emulators are fits, not theories. No accuracy claim can be made for states outside
the sampled envelope. Examples of the behaviour of each model on extrapolation are reported
in Section~\ref{sec:results}. 

\subsection{Data preparation}
\label{sec:prep}

Three operations are applied to prepare \textsc{MELTS} and \textsc{HeFESTo} outputs for training.

\subsubsection{Consistency filtering} Rows are removed where the reference code's output
is internally inconsistent, e.g. where the composition reconstructed from the
reported assemblage disagrees with the stated bulk composition beyond tolerance. The fractions
actually removed are reported in Table~\ref{tab:filters}.

\subsubsection{Rare-phase removal} Phases occurring too infrequently in the training set
to be effectively learned are removed from the component basis entirely to simplify the model.
This operates on entire rows; assemblages with rare phases are discarded
for containing a rare phase at any low modal proportion. It is presumed that this operation is minimally damaging, because very few assemblages are affected by design. 

\subsubsection{Rare-phase augmentation} The \textsc{MELTS} 1.2.0 emulator's training dataset upsampled phases containing rare phases like calcite that were especially desired so that they would surpass the removal threshold and be tracked by the emulator. To upsample, rows containing the rare phase would be duplicated, but with perturbed abundances of its phases, which shift the bulk composition for the same assemblage. This operation was only applied to the isothermal models to avoid invalidating the extensive entropy value natively reported by \textsc{MELTS}. For all models, modal imbalance is addressed through per-phase loss weighting (Section~\ref{sec:phaseweights}) during training, which incentivizes the emulator to learn rare phases that otherwise would have little affect on the loss function.

% -- TABLE PLACEHOLDER ------------------------------------------------------- .... Maybe.... uhh maybe not.
% \begin{table}[t] Rows entering and surviving each filter, per released model. 
% \label{tab:filters} \end{table}
% ----------------------------------------------------------------------------

\subsubsection{Partitioning}
\label{sec:split}

To monitor model overfitting during training, a validation dataset distinct from the training set is used. Training and validation datasets are generated with the same pipelines, but are separated \emph{by unique bulk
composition}. Bulk compositions were partitioned between training and validation datasets and are treated independently as separate datasets throughout the processing pipeline. 
This ensures that the validation dataset tracks the same fundamental structure as the training dataset while avoiding data leakage.
Successive states along one path are strongly autocorrelated, and a row-wise random
split would place near-duplicates on both sides of the partition. The resulting
held-out error would measure interpolation between adjacent states of a rock the
model had already seen, rather than generalization to a rock it had not, and would be
optimistic by a wide and unquantifiable margin. Splitting by composition, before the reference thermodynamic
models were applied, makes the validation dataset's error track overfitting to the training dataset more faithfully. 
This overfitting metric is instrumental to develop an emulator that can be effectively applied to new data. 

Compositions were assigned in the ratio 80:20 between training and validation. 
The validation partition is used every epoch for model selection, to tune adaptive regularization, and for early stopping. 
Unless otherwise stated, all performance metric presented here in Section~\ref{sec:results} are for the validation datasets; the emulator's performance on the training datasets is not reported. 



\subsection{Training}\label{ch4-training}
-continuous model types: Affinity aware and affinity naive
-figure: affinity projection method dndp dndT (plot of real, rescaled affinity vs projection) 
-figure: softplus thermal annealing activation for different a, saturation activation (weighted sigmoid, on the same scale?)

-figure: learning rate schedule + losses, shared epoch x axis

\subsection{Material Properties from Phase Equilibria}\label{ch4-material-props}
-vectorized EOS. Refer to sources
-figure: 1:1 plots of material properties (density, bulk modulus, thermal expansivity, heat capacity) for MELTS and HeFESTo

\section{Results}\label{ch4-results}

-Quality Tables for MELTS and HeFESTo (chemical errors, mole fractions, material properties, deployment tests)
-Figures
    - Phase sections P|S for HeFESTo,  T|P for MELTS
    - Phase diagrams for HeFESTo, MELTS in these same spaces
    - Isentrope properties for HeFESTo
    - Density curves by phase for MELTS?
    - LLD's for MELTS,
    - computational speed comparisons

\section{Discussion}\label{ch4-discussion}
- Compare to other surrogate models in geodynamics (research and citations)
- Discuss model limitations and future work 
- Discuss affinity aware vs naive performance
- use cases: compositionally aware seismic inversion (Mars profiles)
- transition zone search: depth, strength, width of transition zone as a function of composition, plot colored ternaries for depth, strength, width of transition zone